\documentclass{article}
\usepackage{amsmath,amssymb}
\begin{document}

Write \(\varepsilon^j=R^j-u_h^j\), and let \(\widetilde R\) and \(u_h\) be the continuous, piecewise affine interpolants of their nodal values. Set \(\varepsilon(t)=\widetilde R(t)-u_h(t)\). The elliptic-reconstruction bounds give
\[
\|\varepsilon^0\|_{\infty,\Omega}\leq\eta_{\mathrm{ell}}^0,
\qquad
\|\delta_t\varepsilon^j\|_{\infty,\Omega}
\leq\eta_{\mathrm{ell},\delta_t}^j.
\]
For \(j\geq1\), the discrete equation and the definition of \(\psi_h^j\) imply \(\psi_h^j=-\delta_tu_h^j\). Moreover, because \(\mathcal M\) is time-independent and linear, the defining elliptic problems yield, on \(I_j\),
\[
\mathcal M\widetilde R(s)
=f^j+\psi_h^j+(s-t_j)\delta_t(f^j+\psi_h^j),
\qquad \partial_s\widetilde R(s)=\delta_tR^j.
\]
Consequently, for \(w=u-\widetilde R\),
\[
\mathcal Kw(s)
=f(s)-f^j+(t_j-s)\delta_t(f^j+\psi_h^j)
-\delta_t\varepsilon^j,\qquad s\in I_j.
\]
The function \(w\) belongs to the space required by the Green’s-function representation lemma, and \(w(0)=u_0-u_h^0-\varepsilon^0\). Applying that lemma at \(T\), taking absolute values, and using the Green’s-function \(L^1\) bound therefore gives
\[
\begin{aligned}
\|w(T)\|_{\infty,\Omega}
\leq{}&\phi_{0,\Gamma}(T)
 \bigl(\|u_0-u_h^0\|_{\infty,\Omega}+\eta_{\mathrm{ell}}^0\bigr)\\
&+\sum_{j=1}^M\int_{I_j}\phi_{0,\Gamma}(T-s)
 \bigl\|f(s)-f^j+(t_j-s)\delta_t(f^j+\psi_h^j)\bigr\|_{\infty,\Omega}\,ds\\
&+\sum_{j=1}^M\eta_{\mathrm{ell},\delta_t}^j
 \int_{I_j}\phi_{0,\Gamma}(T-s)\,ds.
\end{aligned}
\]
Finally, \(u_h(T)=u_h^M\), \(\widetilde R(T)=R^M\), and \(\|R^M-u_h^M\|_{\infty,\Omega}\leq\eta_{\mathrm{ell}}^M\). Thus the desired a posteriori upper bound is
\[
\boxed{\begin{aligned}
\|u(T)-u_h(T)\|_{\infty,\Omega}
\leq{}&\eta_{\mathrm{ell}}^M
+\phi_{0,\Gamma}(T)
 \bigl(\|u_0-u_h^0\|_{\infty,\Omega}+\eta_{\mathrm{ell}}^0\bigr)\\
&+\sum_{j=1}^M\left[
 \int_{I_j}\phi_{0,\Gamma}(T-s)
 \bigl\|f(s)-f^j+(t_j-s)\delta_t(f^j+\psi_h^j)\bigr\|_{\infty,\Omega}\,ds
 +\eta_{\mathrm{ell},\delta_t}^j
 \int_{I_j}\phi_{0,\Gamma}(T-s)\,ds\right].
\end{aligned}}
\]
Every term on the right is specified by the data, the discrete solution, and the assumed elliptic estimators. The estimate is understood whenever its displayed maximum-norm data terms are finite.

\end{document}